%%%%%%%%%%%%%%%%%%%%% chapter-02.tex %%%%%%%%%%%%%%%%%%%%%%%%%%%%%%%%%
%
% Capítulo 2: Teoría Formal de Autómatas Celulares y Espacios Discretos
%
%%%%%%%%%%%%%%%%%%%%%%%%%%%%%%%%%%%%%%%%%%%%%%%%%%%%%%%%%%%%%%%%

\chapter{Teoría formal de autómatas celulares y espacios discretos}
\label{cap:teoria_formal}

\abstract*{Este capítulo establece el marco matemático fundamental sobre el cual se sustenta la dinámica de los autómatas celulares. A partir de conceptos de la teoría de conjuntos, topología discreta y álgebra geométrica, se construyen las definiciones axiomáticas del espacio celular, los alfabetos finitos, las vecindades y las funciones de transición. Posteriormente, se abordan las propiedades topológicas que permiten el modelado de la computación sobre arquitecturas no convencionales, dotando de formalismo a las simulaciones de sistemas complejos.}

\abstract{Este capítulo establece el marco matemático fundamental sobre el cual se sustenta la dinámica de los autómatas celulares. A partir de conceptos de la teoría de conjuntos, topología discreta y álgebra geométrica, se construyen las definiciones axiomáticas del espacio celular, los alfabetos finitos, las vecindades y las funciones de transición. Posteriormente, se abordan las propiedades topológicas que permiten el modelado de la computación sobre arquitecturas no convencionales, dotando de formalismo a las simulaciones de sistemas complejos.}

\section{Fundamentación conjuntista y topológica del espacio celular}
\label{sec:fundamentacion}

Para formalizar el estudio de la computación en sistemas dinámicos discretos, es imperativo establecer la estructura geométrica base sobre la cual operará el sistema. El análisis del comportamiento computacional y de la emergencia de partículas requiere que el medio físico (o su abstracción matemática) esté rigurosamente delimitado bajo estrictos criterios conjuntistas.


\subsection{Definición del espacio y la unidad celular}
Se define el espacio celular como una retícula geométrica infinita y multidimensional, denotada matemáticamente por el conjunto $\mathbb{Z}^d$, donde $\mathbb{Z}$ representa el conjunto de los números enteros y $d \in \mathbb{N}$ determina la dimensionalidad del espacio analizado \cite{Kari2005}.

Dentro de esta retícula, se define formalmente el concepto de \textbf{célula} como la unidad mínima indivisible del modelo de autómata celular. Cada célula posee una posición vectorial discreta e inmutable en el espacio. Sea $c$ una célula cualquiera en el espacio $d$-dimensional, su ubicación espacial exacta está dada por un vector $\vec{v} \in \mathbb{Z}^d$, tal que:
\begin{equation}
    \vec{v} = (x_1, x_2, \dots, x_d) \quad \text{donde} \quad x_i \in \mathbb{Z}
\end{equation}

Para los propósitos de modelado del presente trabajo, se hará especial énfasis en el análisis del caso bidimensional ($d = 2$). En un espacio bidimensional ortogonal $\mathbb{Z}^2$, la posición de cada célula se simplifica a un par ordenado de coordenadas discretas $(i, j)$. 

\subsection{Definición axiomática de autómata celular}
Un autómata celular (AC) es un sistema dinámico y discreto en el espacio, el tiempo y el estado. Es un modelo que evoluciona mediante pasos de tiempo sincrónicos donde cada elemento evalúa su estado futuro basándose en su entorno local \cite{Wolfram2002}. Matemáticamente, se define como una cuádrupla funcional:
\begin{equation}
    AC = (\mathcal{L}, \Sigma, \mathcal{N}, \delta)
\end{equation}
donde:
\begin{itemize}
    \item $\mathcal{L}$ es el espacio subyacente, constituido por la retícula de dimensión finita (comúnmente $\mathbb{Z}^d$).
    \item $\Sigma$ es el alfabeto o conjunto finito de estados elementales.
    \item $\mathcal{N}$ es la vecindad o subconjunto finito de desplazamientos vectoriales adyacentes a la célula analizada.
    \item $\delta$ es la función de transición local que proyecta el estado actual de la vecindad al estado futuro de la célula central.
\end{itemize}

\subsection{Alfabetos finitos y el espacio de estados}
La dinámica de la información en un autómata celular se manifiesta a través de la alteración sistemática de los valores asignados a cada coordenada espacial. Se define formalmente un conjunto finito de estados discretos, denominado alfabeto y denotado por el símbolo $\Sigma$. Este conjunto contiene íntegramente todos los posibles estados mutuamente excluyentes que una célula $c \in \mathbb{Z}^d$ puede adoptar en un instante de tiempo arbitrario $t$.

Por definición, el alfabeto cumple con la propiedad de cardinalidad finita $|\Sigma| = k$, para todo $k \in \mathbb{N}$ tal que $k \geq 2$. En su forma más reducida y elemental, particularmente en los dominios de la computación lógica clásica y la electrónica digital, el alfabeto empleado es el binario, donde $\Sigma = \{0, 1\}$. 

Adicionalmente a la cardinalidad, se establece la existencia obligatoria de un estado quiescente o estado de reposo absoluto $q_0 \in \Sigma$. El estado quiescente posee la característica de equilibrio local: si una célula arbitraria y la totalidad de sus células adyacentes determinadas por la vecindad $\mathcal{N}$ se encuentran concurrentemente en el estado $q_0$ en el tiempo $t$, la célula central evaluada permanecerá forzosamente en el estado $q_0$ para el paso temporal siguiente $t+1$.

\subsection{Formalización métrica de vecindades acotadas}
El flujo y procesamiento de la información en un autómata celular están estrictamente limitados y supeditados a interacciones puramente locales. Para formalizar mecánicamente este principio de localidad, se definen matemáticamente los subconjuntos finitos adyacentes que influyen sobre la función de transición de una célula central.

Una vecindad genérica $\mathcal{N}$ de tamaño $n$ se define como un conjunto finito y ordenado de vectores de desplazamiento relativo desde el origen:
\begin{equation}
    \mathcal{N} = \{\vec{v}_1, \vec{v}_2, \dots, \vec{v}_n\} \quad \text{donde} \quad \vec{v}_i \in \mathbb{Z}^d
\end{equation}
En consecuencia, para una célula central ubicada en la posición absoluta $\vec{c} \in \mathbb{Z}^d$, el conjunto coordenado de todas sus células vecinas se obtiene mediante la suma vectorial homóloga:
\begin{equation}
    V(\vec{c}) = \{\vec{c} + \vec{v}_i \mid \vec{v}_i \in \mathcal{N}\}
\end{equation}

Dentro de la retícula estándar bidimensional $\mathbb{Z}^2$, las configuraciones topológicas más aplicadas se delimitan utilizando métricas de distancia discreta \cite{Adamatzky1994}. La vecindad de Moore de radio $r=1$ comprende a las 8 células adyacentes (delimitadas por la distancia de Chebyshev infinita normada a 1), mientras que la vecindad de von Neumann incluye exclusivamente a las 4 células ortogonales (determinadas por la distancia de Manhattan igual a 1).

\subsection{El autómata celular hexagonal}


Como punto medular de esta investigación, se introduce una variante topológica de profundo interés analítico: el \textbf{autómata celular hexagonal}. A diferencia de la retícula cartesiana tradicional $\mathbb{Z}^2$ donde cada célula es modelada como un polígono cuadrangular y comparte límites asimétricos (bordes directos frente a vértices diagonales) con sus 8 vecinos de Moore, la teselación hexagonal plana permite dividir uniformemente el espacio mediante polígonos regulares de seis lados.

Esta propiedad geométrica dota al sistema de una simetría isotrópica inmensamente superior, al garantizar que la distancia proyectada desde el centroide de una célula al centroide de todos y cada uno de sus vecinos inmediatos sea matemáticamente idéntica.

En un autómata celular hexagonal, se define el espacio $\mathcal{L}_{hex}$ como una malla isomórfica a una red reticular triangular bidimensional \cite{Basurto2009}. La vecindad básica, equivalente a un radio $r=1$, se compone estrictamente de las seis células adyacentes que comparten de forma colineal una arista con la célula central. 

Formalmente, si se aplica una inmersión de esta red hexagonal sobre un sistema de coordenadas oblicuas bidimensionales $(x, y)$, la vecindad canónica hexagonal $\mathcal{N}_{hex}$ se modela axiomáticamente mediante el conjunto de los seis vectores de desplazamiento relativo:
\begin{equation}
    \mathcal{N}_{hex} = \{ (1,0), (0,1), (-1,1), (-1,0), (0,-1), (1,-1) \}
\end{equation}

El uso del sistema de coordenadas oblicuas permite incrustar la topología de base 6 dentro de los límites matemáticos y computacionales de arreglos matriciales discretos. Esta estructura topológica es fundamental en la disciplina de simulación no convencional, debido a que su uniformidad reduce los sesgos angulares (anisotropía) durante la propagación dinámica de información y modelado de trayectorias de partículas complejas.

% TODO: Insertar diagrama comparativo entre la retícula de cuadrícula (con vecindad de Moore/von Neumann) y la retícula hexagonal. Se debe resaltar la igualdad de distancias desde la célula central hacia sus vecinos en la proyección hexagonal.
\begin{figure}[h]
\centering
% \includegraphics[width=0.8\textwidth]{figuras/comparativa_reticulas.png}
\caption{Comparativa geométrica topológica entre el espacio de cuadrícula $\mathbb{Z}^2$ asimétrica y la teselación celular hexagonal isótropa. Se detallan los vectores de desplazamiento $\vec{v}_i$ que definen la vecindad adyacente $\mathcal{N}_{hex}$.}
\label{fig:topologia_hexagonal}
\end{figure}


\section{Funciones de transición y evolución espacio-temporal}
\label{sec:funciones_transicion}

Una vez establecida la base estructural del espacio celular en la sección anterior, es necesario formalizar el mecanismo que rige el cambio de estado de las células a través del tiempo. La transición de una configuración estática a un sistema dinámico se logra mediante la definición de reglas de evolución locales, las cuales determinan el comportamiento emergente del autómata celular en una escala global.

\subsection{La función de transición local}

La dinámica de un autómata celular está gobernada por una función de transición local $\delta$, la cual actúa como un mapeo determinista entre el estado actual de la vecindad de una célula y su estado futuro. Sea $\sigma_c^t \in \Sigma$ el estado de una célula $c$ en el instante de tiempo $t$. La evolución de dicha célula se define formalmente como:

\begin{equation}
    \sigma_c^{t+1} = \delta(\sigma_{v_1}^t, \sigma_{v_2}^t, \dots, \sigma_{v_n}^t)
\end{equation}

donde $(\sigma_{v_1}^t, \dots, \sigma_{v_n}^t)$ representa la configuración de estados de las células contenidas en la vecindad $\mathcal{N}$ en el tiempo $t$. Es fundamental subrayar que la función $\delta$ es idéntica para todas las células de la retícula y permanece invariante ante traslaciones espaciales y temporales \cite{Kari2005}. Esta uniformidad garantiza que las leyes "físicas" del universo digital sean consistentes en todo el dominio $\mathbb{Z}^d$.

\subsection{Configuraciones globales y el mapeo paralelo}

Para describir el estado del sistema en su totalidad, se introduce el concepto de \textbf{configuración global} $\mathcal{C}$, la cual se define como una función que asigna un estado del alfabeto a cada célula de la retícula: $\mathcal{C}: \mathbb{Z}^d \to \Sigma$. El conjunto de todas las configuraciones posibles se denomina espacio de configuraciones y se denota por $\Sigma^{\mathbb{Z}^d}$.

La evolución global del autómata celular se formaliza mediante el mapeo de configuración global $G: \Sigma^{\mathbb{Z}^d} \to \Sigma^{\mathbb{Z}^d}$. Este mapeo aplica simultáneamente la función de transición local $\delta$ en cada posición de la retícula, de modo que la transición de una configuración $\mathcal{C}_t$ a $\mathcal{C}_{t+1}$ ocurre de forma síncrona y paralela:

\begin{equation}
    \mathcal{C}_{t+1} = G(\mathcal{C}_t)
\end{equation}

Esta transición de lo estático (la configuración individual) a lo dinámico (la secuencia de configuraciones) permite observar la propagación de información y la formación de estructuras complejas. En el caso de los autómatas celulares complejos, este mapeo global es el que da lugar a la fenomenología de partículas o \textit{gliders}, los cuales pueden interpretarse como entidades computacionales que se desplazan sobre el espacio de estados \cite{Wolfram2002}.

\subsection{Discretización del tiempo y trayectoria evolutiva}

El tiempo en este modelo se considera una variable discreta $t \in \mathbb{N} \cup \{0\}$. La evolución completa del sistema partiendo de una configuración inicial $\mathcal{C}_0$ se describe mediante la órbita o trayectoria:

\begin{equation}
    \mathcal{C}_0, G(\mathcal{C}_0), G(G(\mathcal{C}_0)), \dots, G^t(\mathcal{C}_0)
\end{equation}

En la práctica computacional, estas evoluciones suelen representarse mediante diagramas espacio-temporales, donde una dimensión representa el espacio y otra el tiempo. Para el caso bidimensional ($d=2$), la evolución se visualiza como una sucesión de capas o fotogramas (\textit{frames}), permitiendo el análisis topológico de las colisiones entre partículas, las cuales constituyen la base para la implementación de compuertas lógicas en sistemas complejos.

\begin{figure}[h]
\centering
% TODO: [Diagrama esquemático que muestre la aplicación de la función delta sobre una vecindad para producir un estado en t+1, y cómo esto se extiende al mapeo global G entre dos configuraciones completas]
\caption{Representación del mapeo de estados locales a la evolución de la configuración global del sistema.}
\label{fig:mapeo_global}
\end{figure}

% ----------------------------------------------------------------------
% Bibliografía del Capítulo (Formato IEEE)
% ----------------------------------------------------------------------
\begin{thebibliography}{99.}%

\bibitem{Kari2005}
J. Kari, ``Theory of cellular automata: A survey,'' \textit{Theoretical Computer Science}, vol. 334, no. 1-3, pp. 3--33, 2005.

\bibitem{Wolfram2002}
S. Wolfram, \textit{A New Kind of Science}. Champaign, IL, USA: Wolfram Media, 2002.

\bibitem{Adamatzky1994}
A. Adamatzky, \textit{Identification of Cellular Automata}. Londres, Reino Unido: Taylor \& Francis, 1994.

\bibitem{Basurto2009}
R. Basurto Flores y P. A. León Hernández, ``Computación basada en reacción de partículas en un autómata celular hexagonal,'' Trabajo Terminal No. 20080040, Escuela Superior de Cómputo, Instituto Politécnico Nacional, México, 2009.

\end{thebibliography}